% !TeX program = xelatex
\documentclass[UTF8,fontset=fandol,a4paper,11pt]{ctexart}
\usepackage{amsmath,amssymb,amsthm}
\usepackage[margin=26mm]{geometry}
\usepackage[colorlinks=true,linkcolor=blue,citecolor=blue,urlcolor=blue]{hyperref}
\allowdisplaybreaks[2]
\newtheorem{theorem}{定理}
\newtheorem{remark}[theorem]{注}
\renewcommand{\proofname}{证明}
\newcommand{\R}{\mathbb R}
\newcommand{\BMO}{\mathrm{BMO}}
\newcommand{\supp}{\operatorname{supp}}
\newcommand{\C}{\mathcal C}
\newcommand{\one}{\mathbf 1}
\newcommand{\norm}[1]{\left\lVert#1\right\rVert}

\title{Theorem 5.1 的逆命题与 BMO 刻画}
\author{}
\date{2026 年 11 月 4 日}

\begin{document}
\maketitle

\section{问题与假设}

本文讨论 Sukochev--Yang--Zanin--Zhou \cite{SYZZ} 第 3537 页
Theorem 5.1 的逆命题。其结论是交换子满足一个弱 Ces\`aro
性质 $\mathrm{WCP}(2)$。对 Hilbert 变换或任意一个 Riesz 变换，
该估计能够反过来推出符号属于 $\BMO$；但对仅满足该文
Assumption 1.1 的一般 Calder\'on--Zygmund 算子，逆命题不成立。

记 $|E|$ 为 Lebesgue 测度，$\mu(f)$ 为 $|f|$ 的递减重排，并令
\[
 \C g(t)=\frac1t\int_0^t g(s)\,ds,\qquad t>0.
\]
对局部可积函数 $b$，交换子的符号约定与 \cite{SYZZ} 相同：
\[
 [M_b,T]=M_bT-TM_b.
\]
若测试函数 $f$ 有紧支集，则在其支集之外使用积分表达式
\begin{equation}\label{eq:commutator}
 [M_b,T]f(x)
 =\int_{\R^d}(b(x)-b(y))K(x,y)f(y)\,dy.
\end{equation}
以下证明只在相互分离的立方体上使用这个表达式，
不需要预先假设交换子在 $L^2$ 上有界。

我们将 Theorem 5.1 的结论作为假设：存在统一常数 $A<\infty$，
使得对每个具有有限测度支集的
$0\le f\in L^1(\R^d)\cap L^\infty(\R^d)$，均有
\begin{equation}\label{eq:wcp}
 \lambda\bigl|\{x\in\R^d:|[M_b,T]f(x)|>3\lambda\}\bigr|
 \le A\int_0^{s_f}\C\mu(f)(s)\,ds,
 \qquad \lambda\ge\C\mu(f)(s_f),
\end{equation}
其中 $s_f=|\supp f|$；零函数的情形按平凡情形处理。
这里 $A$ 与 $f$、$\lambda$ 无关。

\section{一般 Calder\'on--Zygmund 算子的障碍}

\cite{SYZZ} 的 Assumption 1.1 只要求核的大小及光滑性上界，
没有要求核非退化。因此它允许 $T=0$。此时对任意
$b\in L^1_{\mathrm{loc}}(\R^d)$，都有 $[M_b,T]=0$，
从而 \eqref{eq:wcp} 自动成立。

但是，取 $b(x)=x_1$，对 $Q_L=[-L,L]^d$ 有 $b_{Q_L}=0$ 以及
\[
 \frac1{|Q_L|}\int_{Q_L}|b(x)-b_{Q_L}|\,dx=\frac L2
 \longrightarrow\infty.
\]
所以 $b\notin\BMO(\R^d)$。这说明一般情形下必须对 $T$
补充适当的非退化条件。

\section{Riesz 变换的逆命题}

固定 $j\in\{1,\ldots,d\}$，设 $R_j$ 是 Riesz 变换，其核为
\[
 K_j(x,y)=\gamma_d\frac{x_j-y_j}{|x-y|^{d+1}},
 \qquad \gamma_d>0.
\]
当 $d=1$ 时，它就是通常归一化的 Hilbert 变换。

\begin{theorem}\label{thm:converse}
设 $b\in L^1_{\mathrm{loc}}(\R^d;\mathbb C)$。
如果 \eqref{eq:wcp} 对 $T=R_j$ 成立，则
$b\in\BMO(\R^d)$，并且
\begin{equation}\label{eq:bmo-bound}
 \norm{b}_{\BMO}
 :=\sup_Q\frac1{|Q|}\int_Q|b(x)-b_Q|\,dx
 \le C_d(1+A),
 \qquad b_Q=\frac1{|Q|}\int_Q b(x)\,dx.
\end{equation}
其中上确界取遍所有边平行于坐标轴的立方体，
$C_d$ 只依赖于维数及 Riesz 变换的固定归一化。
事实上，证明只需要 \eqref{eq:wcp} 对有界可测集的特征函数成立。
\end{theorem}

\begin{proof}
先设 $b$ 为实值函数。

\medskip
\noindent\textbf{第一步：提取特征函数测试估计。}
任取立方体 $Q$ 和可测集 $E\subseteq Q$，令 $f=\one_E$。
因为 $0\le\mu(f)\le1$ 且 $s_f\le|Q|$，所以
\[
 \C\mu(f)(s_f)\le1,
 \qquad
 \int_0^{s_f}\C\mu(f)(s)\,ds\le|Q|.
\]
因此，只要 $\lambda\ge1$，就可以使用 \eqref{eq:wcp}，得到
\begin{equation}\label{eq:indicator}
 \lambda\bigl|\{x\in\R^d:
       |[M_b,R_j]\one_E(x)|>3\lambda\}\bigr|
 \le A|Q|.
\end{equation}
这里特意使用右侧较弱的 $A|Q|$，这样无论将支集理解为
非零集还是闭的本质支集，都可以直接使用该估计。
若 $|E|=0$，相应测试函数为零，无须使用支集参数。

\medskip
\noindent\textbf{第二步：构造核具有固定符号的立方体对。}
任取边长为 $\ell$ 的立方体 $Q$，令
\[
 Q'=Q+3\ell e_j,
\]
其中 $e_j$ 为第 $j$ 个坐标单位向量。显然 $|Q'|=|Q|$。
对 $x\in Q'$、$y\in Q$，有
\[
 x_j-y_j\ge2\ell,
 \qquad |x-y|\le\sqrt{d+15}\,\ell.
\]
于是
\begin{equation}\label{eq:kernel-lower}
 K_j(x,y)\ge\frac{\kappa_d}{|Q|}>0,
 \qquad
 \kappa_d:=\frac{2\gamma_d}{(d+15)^{(d+1)/2}}.
\end{equation}

\medskip
\noindent\textbf{第三步：用中位数控制正部。}
取 $b$ 在 $Q'$ 上的一个中位数 $m$，即
\[
 |\{x\in Q':b(x)>m\}|\le\frac{|Q|}{2},
 \qquad
 |\{x\in Q':b(x)<m\}|\le\frac{|Q|}{2}.
\]
定义
\[
 E_+=\{y\in Q:b(y)\ge m\},
 \qquad F_-=\{x\in Q':b(x)\le m\},
 \qquad
 a_+=\frac1{|Q|}\int_{E_+}(b(y)-m)\,dy.
\]
其中 $a_+<\infty$，且 $|F_-|\ge|Q|/2$。
对几乎处处的 $x\in F_-$，积分中的差值具有相同符号，
由 \eqref{eq:commutator} 和 \eqref{eq:kernel-lower} 得
\begin{align}
 |[M_b,R_j]\one_{E_+}(x)|
 &=\int_{E_+}(b(y)-b(x))K_j(x,y)\,dy \notag\\
 &\ge\frac{\kappa_d}{|Q|}
       \int_{E_+}(b(y)-m)\,dy
 =\kappa_d a_+.
 \label{eq:positive-lower}
\end{align}
如果 $a_+>6/\kappa_d$，取
\[
 \lambda=\frac{\kappa_da_+}{6}>1.
\]
由 \eqref{eq:positive-lower} 可知，几乎处处有
$F_-\subseteq\{|[M_b,R_j]\one_{E_+}|>3\lambda\}$。
应用 \eqref{eq:indicator}，得到
\[
 \frac{\kappa_da_+}{6}\,\frac{|Q|}{2}
 \le\lambda|F_-|\le A|Q|,
\]
从而 $a_+\le12A/\kappa_d$。结合 $a_+\le6/\kappa_d$
的另一种情形，总有
\begin{equation}\label{eq:positive-bound}
 a_+\le\frac{\max\{6,12A\}}{\kappa_d}.
\end{equation}

\medskip
\noindent\textbf{第四步：控制负部并得到平均振荡估计。}
令
\[
 E_-=\{y\in Q:b(y)<m\},
 \qquad F_+=\{x\in Q':b(x)\ge m\},
 \qquad
 a_-=\frac1{|Q|}\int_{E_-}(m-b(y))\,dy.
\]
此时 $|F_+|\ge|Q|/2$，并且对几乎处处的 $x\in F_+$，
\[
 |[M_b,R_j]\one_{E_-}(x)|
 =\int_{E_-}(b(x)-b(y))K_j(x,y)\,dy
 \ge\kappa_da_-.
\]
重复上一步的论证可得
\[
 a_-\le\frac{\max\{6,12A\}}{\kappa_d}.
\]
因此
\[
 \frac1{|Q|}\int_Q|b(y)-m|\,dy
 =a_++a_-
 \le\frac{2\max\{6,12A\}}{\kappa_d}.
\]
虽然 $m$ 是在 $Q'$ 上选择的中位数，但它仍可作为 $Q$ 上
平均振荡的比较常数。由
\[
 |b_Q-m|\le\frac1{|Q|}\int_Q|b(y)-m|\,dy
\]
可得
\[
 \frac1{|Q|}\int_Q|b(y)-b_Q|\,dy
 \le\frac2{|Q|}\int_Q|b(y)-m|\,dy
 \le\frac{4\max\{6,12A\}}{\kappa_d}
 \le\frac{48}{\kappa_d}(1+A).
\]
对所有 $Q$ 取上确界，即得到实值符号的结论。

\medskip
\noindent\textbf{第五步：复值符号。}
设 $b=u+iv$，其中 $u,v$ 为实值函数。因为 Riesz 核为实值，
对实值测试函数 $f$，有
\[
 [M_u,R_j]f=\operatorname{Re}([M_b,R_j]f),
 \qquad
 [M_v,R_j]f=\operatorname{Im}([M_b,R_j]f).
\]
所以 $u$、$v$ 都满足具有同一常数 $A$ 的测试估计。
分别应用实值情形，并利用
\[
 \norm{b}_{\BMO}
 \le\norm{u}_{\BMO}+\norm{v}_{\BMO},
\]
即可得到 \eqref{eq:bmo-bound}。
\end{proof}

\section{适用范围与说明}

\begin{remark}[所需的核条件]
对于实值符号，上述证明实际只使用如下条件：对每个立方体 $Q$，
存在一个与它分离、大小相同的立方体 $Q'$，使得 $K$
在 $Q'\times Q$ 上具有固定实符号，且
\[
 |K(x,y)|\ge\frac{\kappa}{|Q|},\qquad x\in Q',\ y\in Q,
\]
其中 $\kappa>0$ 与 $Q$ 无关。
在这一条件下，同样得到 $\norm{b}_{\BMO}\lesssim\kappa^{-1}(1+A)$。
对于实核，复值符号仍可分别处理实部与虚部。
\end{remark}

\begin{remark}[量词与归一化]
如果假设 \eqref{eq:wcp} 对所有满足 \cite{SYZZ} 中
Assumption 1.1 的算子 $T$ 都成立，那么任选一个 $R_j$
就能推出 $b\in\BMO$。
如果只对某个固定的一般 $T$ 假设该估计，则需要额外的核条件。
此外，Theorem 5.1 中 $\norm{b}_{\BMO}=1$ 是正向定理的归一化；
上述逆命题给出的是 BMO 范数的定量上界，
并不意味着可以反推出其范数恰好等于 $1$。
\end{remark}

\begin{remark}[与端点逆刻画的关系]
Accomazzo \cite{Accomazzo} 研究了交换子的端点 $L\log L$
估计对 BMO 的逆刻画。
本文的证明直接使用 \cite{SYZZ} 的 Theorem 5.1 所给出的条件，
且只需特征函数在 $\lambda\ge1$ 时的水平集估计。
整个论证不通过交换子的 $L^2$ 有界性，因此没有预先使用
待证的 BMO 性质。
\end{remark}

\begin{thebibliography}{9}

\bibitem{SYZZ}
F. Sukochev, F. Yang, D. Zanin, and D. Zhou,
\emph{Optimal distribution estimates for commutators and
Marcinkiewicz multipliers},
Mathematische Annalen \textbf{393} (2025), 3513--3560.
\href{https://doi.org/10.1007/s00208-025-03306-3}
{doi:10.1007/s00208-025-03306-3}.

\bibitem{Accomazzo}
N. Accomazzo,
\emph{A characterization of BMO in terms of endpoint bounds for
commutators of singular integrals},
arXiv:1802.05516 (2018).
\url{https://arxiv.org/abs/1802.05516}.

\end{thebibliography}
\end{document}
